\section{Introduction}

Energy storage can transfer electricity across time, charging in relatively inexpensive hours and discharging when the value of energy is higher. The apparent simplicity of this description conceals an intertemporal control problem. A decision in one hour changes the state of charge (SOC), feasible power, and energy available later. Round-trip losses and degradation reduce gross price-spread value, while the operational meaning of a price row depends on when and how it became available. These features make storage arbitrage a useful test case for systems in which predictive accuracy and decision quality are related but not identical.

This distinction is consequential for electricity-price forecasting. Standard benchmark practice emphasizes MAE, RMSE, probabilistic losses, and careful temporal evaluation \cite{lago2021epf}. Those metrics answer whether a price series was predicted accurately. A storage owner, however, needs to know whether the resulting schedule captured valuable hours without violating the physical envelope. Improving errors in uneventful hours may have little operational effect, whereas reversing the order of a price minimum and maximum can change a schedule materially. Forecast evaluation should therefore be complemented by a common downstream decision problem.

Ukraine provides a particularly relevant setting. The energy system operates under exceptional security and flexibility pressures, and storage is potentially valuable as part of a broader resilience portfolio \cite{iea2024ukraine}. At the same time, an academic prototype must not imply authority to participate in a market. Historical and published hourly prices can support analysis and an operator-facing preview, but market execution requires legal responsibility, credentials, submission protocols, receipts, settlement, and operational monitoring. This paper deliberately addresses the narrower and verifiable question: how should forecast-derived BESS schedules be evaluated and guarded before any execution layer exists?

The primary empirical scope is the Ukrainian day-ahead market. Official DAM rows from JSC Market Operator (OREE) provide the local market source \cite{oree2026dam}. Weather context can be assembled from a source such as Open-Meteo under explicit temporal-availability rules \cite{openmeteo2026archive}. When an official row for a target hour is already published, the row is the price source for the preview; a machine-learning model is not used to re-predict a known price. When a target horizon is not yet published, forecast models may generate price scenarios that are then converted into candidate schedules. This ``official row first'' rule prevents a forecast from being presented as an official price and separates source provenance from model performance.

The central question is descriptive: within the preserved retrospective packet, how do prior-only schedule/value selectors compare with conservative controls when every candidate is evaluated in the same physical and economic contour? Schedule/Value Learner V2+ expands a deterministic schedule library and selects among candidates using evidence available before each anchor. It does not optimize a differentiable market-clearing model. Its role is to expose how schedule-family choice changes value under a common evaluator.

This framing places the work near predict-then-optimize and decision-focused learning (DFL). Smart Predict-then-Optimize formalizes the mismatch between prediction loss and downstream objective quality \cite{elmachtoub2022spo}, while recent DFL surveys describe differentiable, surrogate, and gradient-free approaches for learning decisions under constraints \cite{mandi2024dfl}. Storage-specific work has shown why price learning should account for arbitrage outcomes \cite{sang2022essdfl}. The implementation in this paper is intentionally narrower than full differentiable DFL: the optimizer and evaluator are fixed, and the selector uses prior schedule-value evidence rather than back-propagating through a complete predict-then-bid stack. This limitation is part of the claim, not a detail hidden after the results.

The paper makes four bounded contributions.

\begin{enumerate}
\item It defines a reproducible evidence contour in which raw forecasts, calibrated forecasts, heuristic candidates, and learned selectors are converted into feasible hourly schedules and evaluated against the same realized-price oracle.
\item It reports raw and calibrated NBEATSx as co-primary sensitivity sources over 18 dates and five configured BESS profiles, with four adjacent evaluation windows, overlapping/expanding prior histories, capacity-normalized reporting, and no claim of independent replication.
\item It provides a claim-to-artifact and model-lineage audit that separates the main V2+ comparison from a random-forest exact-mirror diagnostic, Hugging Face transformer diagnostics, and a distinct read-model readiness packet.
\item It integrates source provenance and explicit claim boundaries into the system architecture, separating published DAM data, forecast scenarios, IDM read-model support, and any future execution context.
\end{enumerate}

The distinction between an empirical estimate and an engineering artifact is explicit. Regret is measured under a fixed evaluator and time-ordered replay. Dagster materializations, strict contracts, an API, and a dashboard make the evidence inspectable, but a functioning interface is not evidence of trading performance. The repository records manifests, run identifiers, hashes, and result tables \cite{fefelov2026repo}.

The remainder reviews forecast, optimization, DFL, and sequence-model context; formulates the guarded pipeline; describes the retrospective protocol; presents co-primary source sensitivities and smoke diagnostics; explains the operator-preview surface; and closes with limitations and reproducibility statements.
